\bookchapter{Глава 30: Трансляция экрана, презентации и автомобильная интеграция}

\chapterintro{Экран телефона даже при больших размерах — всего 7 дюймов, но ваш телевизор — 55 дюймов, а проектор — 100 дюймов. «Перебросив» изображение с телефона на большой экран, вы получаете совершенно иной уровень восприятия. Эта глава подробно объясняет три частых сценария — трансляцию экрана, презентации и автомобильную интеграцию, — помогая расширить возможности телефона на большие экраны.}

\booksection{I. Три лагеря беспроводной трансляции экрана — не путайте}

Беспроводная трансляция выглядит как «изображение телефона на телевизоре», но базовые технологии бывают трёх видов, и они не взаимозаменяемы:

\begin{center}
\begin{tabularx}{\linewidth}{|X|X|X|}
\hline
\textbf{Технология} & \textbf{Кто использует} & \textbf{Особенности} \tabularnewline
\hline
AirPlay & iPhone/iPad/Mac & Экосистема Apple, самый плавный опыт. Требуется Apple TV или телевизор с поддержкой AirPlay \tabularnewline
Miracast & Android-смартфоны & Прямое подключение Wi-Fi Direct, маршрутизатор не нужен. На Android-смартфонах «беспроводная трансляция экрана» по умолчанию использует именно его \tabularnewline
DLNA & Все платформы & Передаёт только медиафайлы (видео/музыка/фото), не дублирует экран. Подходит для передачи видео на телевизор для воспроизведения \tabularnewline
\hline
\end{tabularx}
\end{center}

\booksubsection{Способ применения на разных платформах:}

\begin{itemize}
  \item \textbf{iPhone}: Пункт управления > Повтор экрана > выберите ваш телевизор/Apple TV. Если телевизор не поддерживает AirPlay, можно купить Apple TV или ТВ-приставку с поддержкой AirPlay.
  \item \textbf{Android}: Потяните шторку уведомлений > найдите кнопку «Беспроводная трансляция»/«Умная трансляция»/«Многоэкранное взаимодействие» > выберите телевизор. Телефон и телевизор должны быть подключены к одной сети Wi-Fi.
  \item \textbf{Универсальное решение}: Если телевизор слишком старый и не поддерживает беспроводную трансляцию, купите Chromecast или приставку Xiaomi/Huawei, подключите к HDMI-порту телевизора — и можно принимать трансляцию экрана.
\end{itemize}

\booksection{II. Проводная трансляция экрана — стабильнее и с нулевой задержкой}

Беспроводная трансляция удобна, но имеет задержку — при играх и презентациях опыт неудовлетворительный. Проводная трансляция решает эту проблему:

\begin{itemize}
  \item \textbf{iPhone}: Требуется переходник Lightning/USB-C на HDMI (около 100–200 юаней). После подключения изображение с телефона напрямую выводится на HDMI-порт телевизора/проектора.
  \item \textbf{Android}: Некоторые смартфоны поддерживают прямое подключение USB-C-HDMI (требуется поддержка DP Alt Mode, есть у большинства флагманов). Достаточно купить кабель USB-C to HDMI.
  \item \textbf{Сценарий переговорной}: Если вам часто приходится проводить презентации с телефона, рекомендуется всегда носить в сумке переходник на HDMI. В переговорных есть телевизоры/проекторы с HDMI — подключил и можно работать, это надёжнее, чем возиться с беспроводной трансляцией.
\end{itemize}

\booksection{III. Телефон превращается в компьютер — режим рабочего стола}

Некоторые Android-смартфоны поддерживают «режим рабочего стола» — при подключении монитора телефон выводит интерфейс, похожий на рабочий стол компьютера, управлять которым можно мышью и клавиатурой:

\begin{itemize}
  \item \textbf{Samsung DeX}: Фирменная функция флагманов Samsung. После подключения монитора появляется интерфейс рабочего стола, можно работать в нескольких окнах — опыт ближе всего к настоящему компьютеру. С Bluetooth-клавиатурой и мышью в командировке можно обойтись без ноутбука.
  \item \textbf{Режим рабочего стола Huawei}: У серий Huawei Mate/Pura при подключении монитора тоже есть похожая функция. Поддерживает многооконность, файловый менеджер и офисные приложения.
  \item \textbf{Motorola Ready For}: Флагманы Moto тоже поддерживают, функции аналогичны DeX.
  \item \textbf{Подходящие сценарии}: Лёгкая офисная работа (ответ на письма, правка документов), презентация PPT, просмотр видео на большом экране. Не подходит для интенсивной работы в Excel или профессионального дизайна.
\end{itemize}

\booksection{IV. Автомобильная интеграция — правильный способ пользоваться телефоном за рулём}

Управлять телефоном во время движения незаконно, но навигация, музыка и приём звонков — насущная необходимость. Системы автомобильной интеграции позволяют использовать ключевые функции телефона на экране автомобильной мультимедиа, а сам телефон можно отложить в сторону и не трогать:

\begin{itemize}
  \item \textbf{Apple CarPlay}: Первый выбор для пользователей iPhone. Через USB или беспроводное подключение к мультимедиа на экране появляется упрощённый интерфейс iPhone — навигация (AutoNavi/Baidu Maps), музыка, телефон, сообщения, подкасты. Голосовое управление Siri позволяет не прикасаться к экрану.
  \item \textbf{Android Auto}: Аналогичное решение для пользователей Android. Функции схожи с CarPlay, поддерживает Google Maps, Spotify и др. В Китае из-за ограничений сервисов Google часть функций недоступна.
  \item \textbf{Huawei HiCar}: Собственное решение Huawei с высокой долей поддержки у автомобилей китайских марок. Помимо навигации и музыки, есть видеозвонки Huawei, совмещение с камерой заднего вида и другие специфические функции.
  \item \textbf{Baidu CarLife}: Кроссплатформенное решение, работает и на iOS, и на Android. Функций меньше, но совместимость шире — многие старые модели автомобилей поддерживают только CarLife.
\end{itemize}

\begin{tcbwarn}{Напоминание о безопасности}
При любом варианте автомобильной интеграции за рулём следует по возможности использовать голосовое управление. Заранее задавайте пункт назначения в навигаторе, заранее выбирайте плейлист — не управляйте экраном мультимедиа во время движения.

\end{tcbwarn}

\booksection{V. Трансляция экрана в переговорной на практике — без суеты}

Совещания — самый частый сценарий трансляции экрана и самый проблемный. Ниже — несколько практических советов:

\begin{itemize}
  \item \textbf{Приходите заранее и проверяйте}: Никогда не подключайте трансляцию в начале совещания. Приходите на 10 минут раньше и проверьте способ подключения. Вероятность сбоя Wi-Fi-трансляции гораздо выше, чем вы думаете.
  \item \textbf{Держите проводной вариант как план Б}: Носите в сумке переходник на HDMI. Когда беспроводная трансляция не подключается, провод — спасательный круг.
  \item \textbf{Закройте предпросмотр уведомлений}: Во время трансляции экрана ваши сообщения WeChat будут отображаться на большом экране переговорной. Перед трансляцией закройте «предпросмотр содержимого уведомлений на экране блокировки» в настройках или включите «режим презентации» (поддерживается на некоторых телефонах: автоматически скрывает уведомления во время трансляции).
  \item \textbf{«Демонстрируйте документ», а не «управляйте телефоном»}: Если нужно показать PPT, откройте его прямо в телефоне через WPS в полноэкранном режиме, а не транслируйте весь рабочий стол. Это профессиональнее и плавнее.
\end{itemize}

\begin{tcbpractice}{Практика: выполните трансляцию экрана с телефона на телевизор}
\begin{enumerate}
  \item Убедитесь, что телефон и телевизор подключены к одной сети Wi-Fi.
  \item Откройте на телефоне фото или видео.
  \item Попробуйте беспроводную трансляцию (на iPhone — «Повтор экрана», на Android — кнопка беспроводной трансляции).
  \item Если не получилось, проверьте, поддерживает ли телевизор нужный протокол, или попробуйте передачу через DLNA (кнопка «трансляция» в приложении «Галерея»/«Видео»).
  \item После успешной трансляции попробуйте выполнять на телефоне другие действия и посмотрите, синхронно ли отображается на большом экране.
\end{enumerate}

\end{tcbpractice}

\begin{tcbsummary}{Итоги главы}
\begin{itemize}
  \item Три лагеря беспроводной трансляции: AirPlay (Apple), Miracast (Android), DLNA (кроссплатформенная передача медиа)
  \item Проводная трансляция стабильнее и без задержки, для переговорной обязателен переходник на HDMI
  \item Samsung DeX/режим рабочего стола Huawei превращают телефон в лёгкий офисный компьютер
  \item Автомобильная интеграция: четыре решения — CarPlay/Android Auto/HiCar/CarLife
  \item Железное правило трансляции в переговорной: проверяйте заранее, готовьте план Б, закройте предпросмотр уведомлений
\end{itemize}

\end{tcbsummary}
